% ═══════════════════════════════════════════════════════════════════
%  Préambule — paquets, style, macros
%  Rapport de projet de fin d'études — Yani Concept by Fati
% ═══════════════════════════════════════════════════════════════════

% ── Langue et encodage ───────────────────────────────────────────
\usepackage[utf8]{inputenc}
\usepackage[T1]{fontenc}
\usepackage[french]{babel}
\usepackage{lmodern}
\usepackage{textcomp}

% Babel français met une espace fine avant « : » et transforme les listes.
% On garde ses règles typographiques, mais on lui interdit de redéfinir les
% puces des listes, pour que itemize reste homogène dans tout le document.
\frenchsetup{StandardLists=true}

% ── Mise en page ─────────────────────────────────────────────────
\usepackage[a4paper,top=2.5cm,bottom=2.5cm,left=3cm,right=2.5cm]{geometry}
\usepackage{setspace}
\onehalfspacing
\usepackage{parskip}
\setlength{\parindent}{0pt}

% ── Graphismes et couleurs ───────────────────────────────────────
\usepackage{graphicx}
\usepackage{xcolor}
\usepackage{float}
\usepackage{tikz}

% ── Tableaux ─────────────────────────────────────────────────────
\usepackage{array}
\usepackage{booktabs}
\usepackage{tabularx}
\usepackage{longtable}
\usepackage{multirow}
\usepackage{colortbl}

% ── Divers ───────────────────────────────────────────────────────
\usepackage{amssymb}
\usepackage{enumitem}
\usepackage{fancyhdr}
\usepackage{titlesec}
\usepackage[font=small,labelfont=bf,labelsep=colon]{caption}

% Babel français légende les tableaux « Table ». L'usage académique marocain,
% et le modèle de rapport suivi ici, disent « Tableau ». La redéfinition passe
% par \addto\captionsfrench, sans quoi babel la réécrirait au \begin{document}.
\addto\captionsfrench{%
  \renewcommand{\tablename}{Tableau}%
  \renewcommand{\listtablename}{Liste des tableaux}%
  \renewcommand{\listfigurename}{Liste des figures}%
}
\usepackage{emptypage}
\usepackage{lastpage}
\usepackage{listings}

% hyperref en dernier (il redéfinit beaucoup de commandes)
\usepackage[hidelinks,breaklinks=true]{hyperref}

% ═══════════════════════════════════════════════════════════════════
%  Charte graphique — reprise de la palette réelle de l'application
%  (or #D8A848 sur crème #F9F6F0, cf. mobile/src/theme/colors.ts)
% ═══════════════════════════════════════════════════════════════════
\definecolor{yaniOr}{HTML}{D8A848}
\definecolor{yaniOrClair}{HTML}{E4C15E}
\definecolor{yaniOrProfond}{HTML}{886028}
\definecolor{yaniNoir}{HTML}{1A1712}
\definecolor{yaniCreme}{HTML}{F9F6F0}
\definecolor{yaniGris}{HTML}{6B6252}
\definecolor{yaniGrisClair}{HTML}{EBE3D4}

% ═══════════════════════════════════════════════════════════════════
%  Titres
% ═══════════════════════════════════════════════════════════════════

% Titre de chapitre : un filet doré, le mot « Chapitre N », puis le titre.
\titleformat{\chapter}[display]
  {\normalfont\huge\bfseries\color{yaniNoir}}
  {\filright\normalsize\color{yaniOrProfond}\MakeUppercase{\chaptertitlename\ \thechapter}}
  {8pt}
  {\filright\Large\color{yaniNoir}}
  [\vspace{4pt}{\color{yaniOr}\titlerule[1.5pt]}]

\titlespacing*{\chapter}{0pt}{0pt}{28pt}

% Chapitre NON numéroté : dédicace, remerciements, résumé, listes, table des
% matières, introduction générale, conclusion, webographie. Titre centré,
% souligné du même filet doré que les chapitres numérotés.
%
% C'est titlesec qui doit porter ce style, et non un titre composé à la main
% avant l'appel : \listoffigures, \listoftables et \tableofcontents ouvrent
% eux-mêmes une page par \chapter*, si bien qu'un titre écrit juste avant se
% retrouverait seul sur la page précédente.
\titleformat{name=\chapter,numberless}[display]
  {\normalfont\bfseries\color{yaniNoir}}
  {}
  {0pt}
  {\centering\LARGE}
  [\vspace{6pt}\centering{\color{yaniOr}\rule{0.35\textwidth}{1.5pt}}]

\titlespacing*{name=\chapter,numberless}{0pt}{0pt}{30pt}

\titleformat{\section}
  {\normalfont\large\bfseries\color{yaniOrProfond}}{\thesection}{0.8em}{}
\titleformat{\subsection}
  {\normalfont\normalsize\bfseries\color{yaniNoir}}{\thesubsection}{0.7em}{}
\titleformat{\subsubsection}
  {\normalfont\normalsize\itshape\color{yaniNoir}}{\thesubsubsection}{0.7em}{}

\setcounter{secnumdepth}{3}
\setcounter{tocdepth}{2}

% ═══════════════════════════════════════════════════════════════════
%  En-têtes et pieds de page
% ═══════════════════════════════════════════════════════════════════
\pagestyle{fancy}
\fancyhf{}
\renewcommand{\headrulewidth}{0.4pt}
\renewcommand{\footrulewidth}{0pt}
\fancyhead[L]{\footnotesize\itshape\color{yaniGris}\nouvelleTitreCourt}
\fancyhead[R]{\footnotesize\itshape\color{yaniGris}\leftmark}
\fancyfoot[C]{\thepage}

% `plain` (première page de chapitre) garde le numéro, sans en-tête.
\fancypagestyle{plain}{%
  \fancyhf{}%
  \renewcommand{\headrulewidth}{0pt}%
  \fancyfoot[C]{\thepage}%
}

% \leftmark sans le « CHAPITRE 1. » automatique
\renewcommand{\chaptermark}[1]{\markboth{#1}{}}

% ═══════════════════════════════════════════════════════════════════
%  Numérotation des figures et tableaux par chapitre (Figure 1.1, …)
% ═══════════════════════════════════════════════════════════════════
\counterwithin{figure}{chapter}
\counterwithin{table}{chapter}

% ═══════════════════════════════════════════════════════════════════
%  Figures : image réelle si le fichier existe, cadre d'attente sinon
% ═══════════════════════════════════════════════════════════════════
%
%  C'est le mécanisme qui permet de rédiger et de compiler le rapport AVANT
%  d'avoir produit les captures d'écran et les diagrammes UML. Chaque appel
%  cherche le fichier dans `figures/` :
%
%   • le fichier est là  → l'image est insérée à la largeur demandée ;
%   • le fichier manque  → un cadre gris s'affiche à sa place, portant le nom
%                          de fichier attendu et la légende.
%
%  Il n'y a donc RIEN à modifier dans le texte le jour où les images arrivent :
%  il suffit de déposer `figures/<nom>.png` et de recompiler.
%
%  Usage :  \figProjet[largeur]{nom-du-fichier}{Légende}{cle-de-label}
%           largeur  : fraction de \textwidth (0.85 par défaut)
%           renvoi   : \figref{cle-de-label}
%
\newcommand{\cadreAttente}[2]{%
  \setlength{\fboxrule}{1pt}%
  \setlength{\fboxsep}{0pt}%
  \fcolorbox{yaniOr}{yaniCreme}{%
    \begin{minipage}[c][4.2cm][c]{0.85\textwidth}
      \centering
      {\color{yaniOrProfond}\Large$\square$}\\[6pt]
      {\color{yaniGris}\small\textsc{figure à insérer}}\\[8pt]
      {\color{yaniNoir}\small\textbf{#2}}\\[8pt]
      {\color{yaniGris}\footnotesize\ttfamily figures/#1}
    \end{minipage}%
  }%
}

\newcommand{\figProjet}[4][0.85]{%
  \begin{figure}[H]
    \centering
    \IfFileExists{figures/#2}%
      {\includegraphics[width=#1\textwidth,keepaspectratio]{figures/#2}}%
      {\cadreAttente{#2}{#3}}%
    \caption{#3}
    \label{fig:#4}
  \end{figure}%
}

% Deux figures côte à côte (captures d'écran mobile, souvent étroites).
\newcommand{\figDouble}[6][0.42]{%
  \begin{figure}[H]
    \centering
    \begin{minipage}[t]{0.48\textwidth}
      \centering
      \IfFileExists{figures/#2}%
        {\includegraphics[width=#1\textwidth,keepaspectratio]{figures/#2}}%
        {\cadreAttenteEtroit{#2}{#3}}%
    \end{minipage}\hfill
    \begin{minipage}[t]{0.48\textwidth}
      \centering
      \IfFileExists{figures/#4}%
        {\includegraphics[width=#1\textwidth,keepaspectratio]{figures/#4}}%
        {\cadreAttenteEtroit{#4}{#5}}%
    \end{minipage}
    \caption{#3 · #5}
    \label{fig:#6}
  \end{figure}%
}

\newcommand{\cadreAttenteEtroit}[2]{%
  \setlength{\fboxrule}{1pt}%
  \setlength{\fboxsep}{0pt}%
  \fcolorbox{yaniOr}{yaniCreme}{%
    \begin{minipage}[c][4.2cm][c]{0.92\textwidth}
      \centering
      {\color{yaniOrProfond}\large$\square$}\\[5pt]
      {\color{yaniGris}\scriptsize\textsc{figure à insérer}}\\[6pt]
      {\color{yaniNoir}\scriptsize\textbf{#2}}\\[6pt]
      {\color{yaniGris}\scriptsize\ttfamily figures/#1}
    \end{minipage}%
  }%
}

\newcommand{\figref}[1]{figure~\ref{fig:#1}}
\newcommand{\tabref}[1]{tableau~\ref{tab:#1}}

% ═══════════════════════════════════════════════════════════════════
%  Tableaux — style commun
% ═══════════════════════════════════════════════════════════════════
\newcolumntype{L}[1]{>{\raggedright\arraybackslash}p{#1}}
\newcolumntype{C}[1]{>{\centering\arraybackslash}p{#1}}
\renewcommand{\arraystretch}{1.25}

% En-tête de tableau : fond doré pâle, texte sombre.
\newcommand{\entete}[1]{\cellcolor{yaniGrisClair}\textbf{#1}}

% ═══════════════════════════════════════════════════════════════════
%  Code source
% ═══════════════════════════════════════════════════════════════════
\lstdefinestyle{yani}{
  basicstyle=\ttfamily\scriptsize,
  keywordstyle=\color{yaniOrProfond}\bfseries,
  commentstyle=\color{yaniGris}\itshape,
  stringstyle=\color{yaniOrProfond},
  numbers=none,
  frame=single,
  rulecolor=\color{yaniGrisClair},
  backgroundcolor=\color{yaniCreme},
  breaklines=true,
  showstringspaces=false,
  columns=fullflexible,
  xleftmargin=6pt,
  framexleftmargin=6pt,
  literate=%
    {é}{{\'e}}1 {è}{{\`e}}1 {ê}{{\^e}}1 {à}{{\`a}}1
    {ç}{{\c c}}1 {ù}{{\`u}}1 {û}{{\^u}}1 {î}{{\^i}}1 {ô}{{\^o}}1
    {É}{{\'E}}1 {—}{{---}}1 {’}{{'}}1 {«}{{\guillemotleft}}1 {»}{{\guillemotright}}1
}
\lstset{style=yani}

% ═══════════════════════════════════════════════════════════════════
%  Pages liminaires
% ═══════════════════════════════════════════════════════════════════
%
%  Ouvre une page non numérotée : titre centré (style titlesec ci-dessus),
%  entrée dans la table des matières, et en-tête courant correspondant.
%  L'ancre hyperref est posée AVANT le titre, pour que le lien de la table des
%  matières arrive en haut de la page et non sous le titre.
\newcommand{\pageLiminaire}[1]{%
  \cleardoublepage
  \phantomsection
  \chapter*{#1}%
  \addcontentsline{toc}{chapter}{#1}%
  \markboth{#1}{}%
}

% ═══════════════════════════════════════════════════════════════════
%  Encadré « point d'attention » — utilisé pour les pièges d'exploitation
% ═══════════════════════════════════════════════════════════════════
\newenvironment{attention}[1][Point d'attention]{%
  \par\vspace{6pt}
  \begin{list}{}{%
    \setlength{\leftmargin}{12pt}\setlength{\rightmargin}{12pt}
    \setlength{\topsep}{0pt}}
  \item[]
  \setlength{\fboxrule}{0pt}%
  {\color{yaniOrProfond}\textbf{#1.}}\ %
}{%
  \end{list}\vspace{6pt}\par
}

% ═══════════════════════════════════════════════════════════════════
%  Raccourcis de nommage
% ═══════════════════════════════════════════════════════════════════
\newcommand{\yani}{\textit{Yani Concept by Fati}}

% Valeur restant à renseigner : affichée en gris italique, elle se repère d'un
% coup d'œil dans le PDF et se coupe normalement en fin de ligne — au
% contraire d'une suite de points de suspension, qui déborde des cellules de
% tableau étroites.
\newcommand{\aCompleter}[1]{\textit{\textcolor{yaniGris}{#1}}}
